\section*{結言}
今回の5件には、AIを企業の既存インフラに載せる動きと、AIが利用者に代わって継続的に作業する動きが共通している。Google CloudへのOpenAIモデル提供は調達・運用の選択肢を広げ、Claude Duetはエージェントを日常的な作業の担い手へ近づける試みといえる。

同時に、PHPの脆弱性発見支援とLiteLLMの供給網侵害は、AIによる防御の可能性と、AIを支えるソフトウェアの安全性という対照的な課題を浮かび上がらせた。成果の評価では、発表者の主張、独立した発見、影響範囲を切り分ける必要がある。

Gemini 4 Argonの3Dデモも含め、今後の注目点は発表段階の機能がどこまで再現可能で、利用者にどのような実益をもたらすかにある。本号は観測窓内のX投稿を入口とした速報であり、詳細な技術評価や一次資料の再調査は週刊版などの工程に委ねる。
